% Part: many-valued-logic
% Chapter: sequent-calculus
% Section: rules-and-proofs

\documentclass[../../../include/open-logic-section]{subfiles}

\begin{document}

\olfileid{mvl}{seq}{rul}

\olsection{Aturan lan \usetoken{P}{derivation}}

Ing ngisor iki, $\Gamma, \Delta, \Pi, \Lambda$ makili
rerangkening winates saka !!{sentence}s.

\begin{defn}[Sekuen]
Sawijining \emph{sekuen $n$ sisi} iku ekspresi awujud
\[
\Gamma_1 \nSequent \dots \nSequent \Gamma_n
\]
ing ngendi saben $\Gamma_i$ iku rerangkening winates
(bisa uga kothong) saka !!{sentence}s ing basa~$\Lang L$.
\end{defn}

\begin{defn}[Sekuen Wiwitan]
Sawijining \emph{sekuen wiwitan $n$ sisi} iku sekuen $n$ sisi
awujud $!A \nSequent \dots \nSequent !A$ kanggo saben
!!{sentence} $!A$ ing basa kasebut.

Yen basa iku ngemot konektif $0$-argumen~$\star$,
yaiku tetapan proposisional, kita uga njupuk sekuen
$\dots \nSequent \star \nSequent \dots$ kang $\star$ mapan
ing panggonan kanggo nilai kabeneran kang cocog karo~$\tf{\star}
\in V$, dene panggonan liyane kothong.
\end{defn}

Kanggo saben konektif ing logika $n$ nilai~$\Log{L}$,
ana aturan logis kanggo saben nilai kabeneran kang bisa
diwenehake dening konektif kasebut ing~$\Log{L}$.
!!^{derivation}s ing kalkulus sekuen $n$ sisi kanggo~$\Log{L}$
iku wit-wit sekuen; sekuen kang paling ndhuwur yaiku sekuen
wiwitan, lan yen sekuen ana ing sangisore siji utawa luwih
sekuen liyane, sekuen iku kudu manut kanthi bener marang
aturan inferensi kanggo konektif ing~$\Log{L}$.

\begin{defn}[Teorema]
!!^a{sentence}~$!A$ iku \emph{teorema} ing logika $n$ nilai~$\Log{L}$
yen ana !!a{derivation} kanggo sekuen $n$ sisi kang ngemot $!A$
ing saben panggonan kang cocog karo nilai kabeneran pinilih
ing~$\Log{L}$. Kita nulis $\Proves[\Log{L}] !A$ yen $!A$
iku teorema lan $\Proves/[\Log{L}] !A$ yen dudu teorema.
\end{defn}

\begin{defn}[!!^{derivability}]
!!^a{sentence}~$!A$ \emph{!!{derivable} saka} himpunan
!!{sentence}s~$\Gamma$ ing logika $n$ nilai~$\Log{L}$,
$\Gamma \Proves[\Log{L}] !A$, yen lan mung yen ana subhimpunan
winates~$\Gamma_0 \subseteq \Gamma$ lan rerangkening $\Gamma_0'$
saka !!{sentence}s ing~$\Gamma_0$ saéngga sekuen iki
nduweni !!a{derivation}:
\[ \Lambda_1 \nSequent \dots \nSequent \Lambda_n \]
ing ngendi $\Lambda_i$ iku $!A$ yen panggonan $i$ cocog karo
nilai kabeneran pinilih, lan $\Gamma_0'$ yen ora. Yen $!A$
ora !!{derivable} saka $\Gamma$, kita nulis
$\Gamma \Proves/[\Log{L}] !A$.
\end{defn}

Umpamane, logika \L ukasiewicz $3$ nilai nduweni
kalkulus sekuen $3$ sisi. Ing sekuen $3$ sisi
$\Gamma \nSequent \Pi \nSequent \Delta$, $\Gamma$ cocog
karo~$\False$, $\Delta$ karo~$\True$, lan $\Pi$ karo~$\Undef$.
Aksiomane $!A \nSequent !A \nSequent !A$. Amarga mung
$\True$ kang dadi nilai pinilih, $\Gamma \Proves[\LogLuk[3]] !A$
yen lan mung yen sekuen $\Gamma \nSequent \Gamma \nSequent !A$
nduweni !!a{derivation}. (Yen $\Undef$ uga dadi nilai pinilih,
kita bakal mbutuhake !!a{derivation} kanggo
$\Gamma \nSequent !A \nSequent !A$.)

\end{document}
